\documentclass[10pt,a4paper]{amsart}
\usepackage[margin=19mm]{geometry}
\usepackage{amsmath,amssymb,mathtools}
\usepackage[scr=rsfs,cal=euler]{mathalfa}
\usepackage{xfrac}
\usepackage[unicode,hidelinks,bookmarks=false]{hyperref}
\newcommand{\ie}{i.e.\ }
\newcommand{\cf}{cf.\ }
\newcommand{\resp}{respectively\ }
\newcommand{\Subsection}{\subsection}
\providecommand{\nobreakdash}{\nobreak\hbox{-}}
\setlength{\emergencystretch}{2em}
\multlinegap0pt
\usepackage{bm}
\usepackage{bbm}
\usepackage{dsfont}
\usepackage{xspace}
\usepackage{cancel}
\usepackage{mathtools}
\usepackage{amsthm}
\makeatletter
\@ifundefined{lemma}{\newtheorem{lemma}{Lemma}}{}
\makeatother
\title[Decoding Quantum LDPC Codes using Collaborative Check Node R]{Decoding Quantum LDPC Codes using Collaborative Check Node Removal}
\author{Mainak Bhattacharyya, Ankur Raina}
\date{Source: arXiv:2501.08036v3 (CC BY 4.0)}
\begin{document}
\maketitle

\noindent\emph{Source and licence.} Decoding Quantum LDPC Codes using Collaborative Check Node Removal. Version arXiv:2501.08036v3 (CC BY 4.0), distributed under the Creative Commons Attribution 4.0 International licence, as recorded in the arXiv licence metadata for this exact version and shown on the abstract page https://arxiv.org/abs/2501.08036v3. Full rights notice: licenses/arxiv-2501.08036-CC-BY-4.0.txt.\par\smallskip

\noindent\textit{Editorial scope.} Counted unit: the Lemma whose source label is \texttt{cor:sep-imp-rem-chk} in the article (source0.tex lines 891--894, source label \texttt{cor:sep-imp-rem-chk}), delivered together with the article's own complete original proof of it (source0.tex lines 895--926, one \texttt{proof} environment of 32 lines). No mathematics is written, repaired or re-derived: statement and proof are the article's own text line for line. Required context retained uncounted: none. Every definition, display and label of those lines is retained unchanged. Omitted from this package and not redistributed: 1--890, 927--1514. Nothing inside a retained excerpt is omitted or altered: every retained body is delivered exactly as the article's lines are written, so no omission receipt of any kind accompanies this package. Citations: the retained excerpts carry no citation command at all, so no bibliography entry of the article is transcribed and no reference is redistributed. Reference adaptation: none was needed and none was made; every reference inside the delivered unit resolves inside it, so no unresolved reference or citation is produced. Printed slips kept literal: no printed word, symbol, hypothesis or number of the counted statement or of its proof was altered. Layout and class: the article's own class is \texttt{revtex4-2} with packages including graphicx, dcolumn, bm, hyperref, amsthm, braket, xcolor, tikz, algorithm2e, subfig; the wrapper is the lane's pinned \texttt{amsart} editorial wrapper at 10pt on A4, which restates the theorem-like environments the retained text needs and adds the article's own macro definitions that the retained text needs (none), with extra wrapper packages (bm, bbm, dsfont, xspace, cancel, mathtools). The delivered numbering is the unit's own. Assets: no retained span contains a figure environment, a graphics-inclusion command, a PDF-page inclusion, a table or a TikZ picture; no image file is copied into this package, the package declares no asset dependency and no image file is redistributed with it. Licence: version arXiv:2501.08036v3 of this article is distributed under the Creative Commons Attribution 4.0 International licence, as recorded in the arXiv licence metadata for this exact version and shown on the abstract page https://arxiv.org/abs/2501.08036v3. A full rights notice is delivered at licenses/arxiv-2501.08036-CC-BY-4.0.txt. Characters: every retained excerpt is pure ASCII, so the delivered engine, XeLaTeX with its default Latin Modern text font, reports no missing character.
\begin{lemma}
    Any erroneous qubit $v_r \in V_i$, where $V_i$ is one of the isomorphic degenerate data qubit subsets of the QTS; can be corrected by \textit{min-sum} BP decoder if the $d_v(d_v-1)$ number of check nodes from the level $2$ of $T(v_r)$, which are limiting the separation of $v_r$; are removed from the computation tree of the decoder.
    \label{cor:sep-imp-rem-chk}
\end{lemma}

\begin{proof}
    To prove the above lemma, we analyze the message received by one of the erroneous data qubits in the QTS.
    The computation tree traversed in root-to-leaf order represents the exact sequence of messages passed by the min-sum BP decoder.
    Therefore, each iteration of the decoder indicates a level in the computation tree.
    % Therefore, we argue that the harmful error beliefs contributed by $c_6$ causes the non-convergence of $v_0$.
    Now assume $v_r \in V_i$, where $V_i$ is one of the degenerate qubit subsets of the QTS and $c_r \in N(V_i)$ is one of the common stabilizer check nodes, which connects two different disjoint qubit subsets of the QTS.
    Also, $T(v_r)$ is the computation tree with $v_r$ as its root.
    We now observe the flow of information in $T(v_r)$ from leaf nodes towards the root node.
    For the \textit{min-sum} BP decoder to be able to correct erroneous $v_r$, correct information must be passed from $c_r$ to root $v_r$ in $T(v_r)$.
    For this to be done, all the information passed to $c_r$ from its descendants must be correct.
    This correct passage of information must satisfy the following equality for messages passed from any child node $v_i$ to the parent $c_r$,
    \begin{align}
        \mathrm{sign}(\lambda_i) = \mathrm{sign}(m_{v_i \rightarrow c_r}); \forall v_i \in N(c_r)\backslash v_r.
        \label{eq:correct-info-sign}
    \end{align}
    Eq. \eqref{eq:correct-info-sign} implies that, as information passed into $c_r$ is correct, the sign of all the messages passed into $c_r$ should be the same as that of the message passed into data qubit nodes $v_i$ from the channel.
    Now, assuming a phenomenological bit flip noise channel, the magnitude of the messages received by the data qubit nodes from the channel is equal.
    Therefore, for $d_v$ regular QLDPC code, the messages sent from any data qubit $v_i$ in the same level of $T(v_r)$ have equal magnitude and follow
    \begin{align}
        |m_{v_i \rightarrow c_j}| &= |\lambda_i| + \sum_{c_{j^\prime}:N(v_i)\backslash c_j}|m_{c_{j^{\prime}}\rightarrow v_i}|.
    \end{align}
    Therefore, if we consider the effect of the stabilizer check nodes, which belong to the common neighborhood of the disjoint subsets, at level $2$ of $T(v_r)$, we directly observe a violation of Eq. \eqref{eq:correct-info-sign}.
    This is because of the following condition
    \begin{align}
        \mathrm{sign}(m_{c_j \rightarrow v_i}) \neq \mathrm{sign}(\lambda_i),
        \label{eq:violate-info-sign}
    \end{align}
    where if $v_r \in V_i$, then $v_i \in V_j$; such that $V_i \cap V_j = \emptyset$.
    At the junction of level $1$ and level $2$, qubit $v_i$ is adjacent to $d_v - 1$ child check nodes, and all such nodes follow Eq. \eqref{eq:violate-info-sign}.
    Therefore, the check nodes at level $2$ of $T(v_r)$, which limit the separation of $v_r$, also pass incorrect information to check node $c_r$ adjacent to $v_r$, which leads to net incorrect information passed to $v_r$ by the \textit{min-sum} BP decoder.
    Therefore, removal of all such $d_v(d_v - 1)$ number of check nodes from level $2$ of $T(v_r)$ ensures correct information passed to the erroneous qubits and, as a result, ensures correctness of the erroneous qubits in the QTS.
\end{proof}

\end{document}
